\documentclass[10pt,a4paper,ragged2e,withhyper]{altacv}


\geometry{left=1.25cm,right=1.25cm,top=1.5cm,bottom=1.5cm,columnsep=1.2cm}

\usepackage{paracol}
\usepackage[hidelinks]{hyperref}

\iftutex

  \setmainfont{Lato}
\else
  \usepackage[default]{lato}
\fi

\definecolor{VividPurple}{HTML}{3E0097}
\definecolor{SlateGrey}{HTML}{2E2E2E}
\definecolor{LightGrey}{HTML}{666666}
% \colorlet{name}{black}
% \colorlet{tagline}{PastelRed}
\colorlet{heading}{VividPurple}
\colorlet{headingrule}{VividPurple}
% \colorlet{subheading}{PastelRed}
\colorlet{accent}{VividPurple}
\colorlet{emphasis}{SlateGrey}
%\colorlet{body}{Black}

\renewcommand{\cvItemMarker}{{\small\textbullet}}


\begin{document}
\name{BEN-HAMID Walid}
\tagline{Ingénieur en Développement Logiciel – Full Stack & Systèmes Web}

\photoR{5cm}{DSC_9057.JPG}
% \photoL{2cm}{Yacht_High,Suitcase_High}
\personalinfo{%
  % Not all of these are required!
  % You can add your own with \printinfo{symbol}{detail}
  \email{wbenhamid4@gmail.com}
   \phone{+212600018820}
  \location{Maroc}
 % \homepage{marissamayr.tumblr.com}
  % \twitter{@marissamayer}
  %\xtwitter{@marissamayer}
  \linkedin{benhamidwalid}

}

\makecvheader

%% Depending on your tastes, you may want to make fonts of itemize environments slightly smaller
\AtBeginEnvironment{itemize}{\small}

%% Set the left/right column width ratio to 6:4.
\columnratio{0.6}

% Start a 2-column paracol. Both the left and right columns will automatically
% break across pages if things get too long.
\begin{paracol}{2}

\cvsection{Expériences Professionnelles et Projets Académiques}
\cvevent{PFE – SonarTec v2 \& Lancement de la plateforme SonarTec-IA}{AXONE SYSTEMS}{Janvier 2025 -- Juillet 2025}{Rabat, Maroc}
\begin{itemize}
  \item Refonte complète de la plateforme \textbf{SonarTec} initiée en PFA, avec une nouvelle architecture logicielle, de nouveaux diagrammes UML et BPMN, et une amélioration du backend avec \textbf{Spring Boot} (remplaçant le code initial généré par \textbf{JHipster}).
  \item Développement d'une interface frontend entièrement repensée avec \textbf{Angular}, pour une meilleure expérience utilisateur.
  \item Intégration du service \textbf{Zoho Recruit} pour importer automatiquement des offres d'emploi, gérer les candidatures et permettre la signature électronique des contrats.
  \item Création de la plateforme complémentaire \textbf{SonarTec-IA} dédiée à l'automa-tisation des entretiens :
  \begin{itemize}
    \item Ajout d'offres et de candidats depuis une interface moderne développée avec \textbf{React}.
    \item Génération dynamique des questions d'entretien via \textbf{OpenAI}.
    \item Enregistrement des réponses vidéos des candidats directement dans la plateforme.
    \item Analyse automatique des performances du candidat (expressions faciales, stress, mouvements) à l'aide de l'IA.
    \item Transcription des réponses et génération d'un rapport d'évaluation intelligent, envoyé aux clients ayant publié l'offre sur \textbf{SonarTec}.
  \end{itemize}
  \item Utilisation de \textbf{Node.js} pour le backend de SonarTec-IA, avec \textbf{Supabase} et \textbf{MySQL} pour la gestion des données.
  \item Collaboration agile avec une équipe interdisciplinaire pour garantir qualité, scalabilité et intégration cohérente entre les deux plateformes.
  \item Gestion de version collaborative avec \textbf{GitHub} et mise en place de workflows CI/CD pour l'automatisation des tests et du déploiement.
\end{itemize}
\cvtag{Spring Boot} \cvtag{Angular} \cvtag{Flowable} \cvtag{Zoho Recruit} \cvtag{React} \cvtag{Node.js} \cvtag{OpenAI} \cvtag{Supabase} \cvtag{MySQL} \cvtag{Vision AI} \cvtag{GitHub} \cvtag{CI/CD}


\divider

\cvevent{PFA – Version 1 : Conception et création de la plateforme \textit{SonarTec}}{AXONE SYSTEMS}{Avril 2024 -- Novembre 2024}{Rabat, Maroc}
\begin{itemize}
\item Conception et modélisation de l'architecture du projet à l'aide de diagrammes UML, y compris des cas d'utilisation, des diagrammes de classes et des workflows BPMN avec \textbf{Flowable}.  
\item Développement du backend en utilisant le framework \textbf{Spring Boot}, avec des outils comme \textbf{JHipster} pour la génération d'applications, et intégration de \textbf{SpringAI} pour les fonctionnalités d'intelligence artificielle.  
\item Création de l'interface utilisateur en \textbf{Angular}, avec une attention particulière à l'ergonomie et à l'expérience utilisateur.  
\item Gestion de la base de données avec \textbf{MySQL}, en assurant une structuration optimale des données pour des performances accrues.  
\item Déploiement d'une plateforme en ligne de recrutement, intégrant des workflows de bout en bout pour simplifier et automatiser les processus RH.  
\item Collaboration avec une équipe multidisciplinaire pour respecter les délais et garantir la qualité du livrable.
\item Utilisation de \textbf{GitLab} pour la gestion de version et la collaboration en équipe, avec mise en place de pipelines CI/CD pour l'automatisation des tests et du déploiement.
\end{itemize}
\cvtag{SpringBoot} \cvtag{Angular} \cvtag{MySQL} \cvtag{SpringAI} \cvtag{GitLab} \cvtag{CI/CD}

\divider




\cvevent{Développement d'une application de gestion des clients}{Java, Spring Boot}{2023}{}
\begin{itemize}
\item Création d'une application web complète pour gérer les profils clients, incluant :
  \begin{itemize}
    \item Gestion des utilisateurs avec rôles et permissions.
    \item Intégration de Spring Security pour l'authentification.
    \item Utilisation de MySQL pour la persistance des données.
  \end{itemize}
\end{itemize}

\divider

\cvevent{Application de gestion des étudiants}{JEE, Hibernate, MySQL}{2024}{}
\begin{itemize}
\item Conception et implémentation d'une plateforme pour la gestion des étudiants :
  \begin{itemize}
    \item Ajout, modification et suppression des informations des étudiants.
    \item Interface utilisateur développée avec JSP et servlets.
    \item Utilisation de Hibernate pour l'accès et la manipulation des données.
  \end{itemize}
\end{itemize}

\divider

\cvevent{Système de gestion des stocks basé sur les automates}{Python, Théorie des Automates}{2024}{}
\begin{itemize}
\item Conception et développement d'un système de gestion des stocks utilisant les automates finis :
  \begin{itemize}
    \item Modélisation des états et transitions pour la gestion des flux d'entrées et de sorties.
    \item Implémentation de la déterminisation, complétion et minimisation des automates.
    \item Visualisation graphique des graphes de transition avec des bibliothèques Python.
    \item Simulation des scénarios réels pour assurer l'optimisation et la robustesse du système.
  \end{itemize}
\end{itemize}



% use ONLY \newpage if you want to force a page break for
% ONLY the currentc column


\switchcolumn




\cvsection{Compétences}

% Frameworks et outils back/front
{\LaTeXraggedright 
\cvtag{SpringBoot} \cvtag{Angular} \cvtag{React} \cvtag{Node.js} \cvtag{Express.js} 
\cvtag{Hibernate/JPA} \cvtag{JHipster} \cvtag{SpringAI} \par}

\divider\smallskip

% Langages de programmation
{\LaTeXraggedright 
\cvtag{C} \cvtag{HTML/CSS} \cvtag{JavaScript} \cvtag{Python} \cvtag{JEE/Java} 
\cvtag{PHP} \cvtag{R} \par}

\divider\smallskip

% Bases de données et Backend as a Service
{\LaTeXraggedright 
\cvtag{MySQL} \cvtag{PostgreSQL} \cvtag{Supabase} \par}

\divider\smallskip

% DevOps et outils de versioning
{\LaTeXraggedright 
\cvtag{Git} \cvtag{GitHub} \cvtag{GitLab} \cvtag{Docker} \cvtag{CI/CD} \par}

\divider\smallskip

% Intelligence Artificielle et Cloud
{\LaTeXraggedright 
\cvtag{GCP} \cvtag{Vertex AI} \cvtag{Speech-to-Text} 
\cvtag{Vision AI} \cvtag{OpenAI} \par}

\divider\smallskip

% Modélisation
{\LaTeXraggedright 
\cvtag{UML} \cvtag{BPMN} \par}

\cvsection{Langues}

\vspace{0.5em}
\begin{tabular}{@{}ll@{}}
\textbf{Français} & Courant \\
\textbf{Anglais} & Intermédiaire \\
\end{tabular}




\cvsection{Formation}

\cvevent{Cycle d'ingénieur en Génie Mathématiques et Informatiques}{Faculté des sciences et techniques Mohammedia}{2022 -- 2025}{Mohammedia}

\divider

\cvevent{Diplôme d'Études Universitaires Scientifiques et Techniques (DEUST), Option MIP}{Faculté des sciences et techniques Mohammedia}{2020 -- 2022}{Mohammedia}

\divider

\cvevent{Baccalauréat Sciences physiques}{Lycée Hassan 2, Benslimane}{2019 -- 2020}{Benslimane}



\cvsection{Centres d'intérêt}

\cvevent{Membre du club de bénévolat}{Rotaract}{2021 -- 2025}{Mohammedia}

\divider

\cvevent{Participation à des cours de soutien en anglais}{Centre Américain}{2023 -- 2024}{Mohammedia}


\cvsection{Loisirs}

{\LaTeXraggedright \cvtag{Football}
\cvtag{Musculation}
\cvtag{Écoute de podcasts}
\cvtag{Voyages et exploration culturelle}
\par}

\end{paracol}

\end{document} 